\newgeometry{top=1cm, left=1cm, right=1cm, bottom=1cm}
\begin{center}
  \includegraphics[height=.6in]{../resources/univ-lyon}%
  % \vrule%
  \includegraphics[height=.6in]{../resources/insa}\\[3pt]%
  \fontsize{14pt}{14pt}\selectfont%
  \MakeUppercase{FOLIO ADMINISTRATIF}\\[3pt]%
  \MakeUppercase{Thèse de l'Université de Lyon opérée au sein de l'Insa Lyon}
\end{center}

\begin{framed}
\noindent %
Nom : \textbf{\MakeUppercase{Brodu}} \hfill Date de soutenance : \textbf{23/06/2016}\\[3pt]
Prénoms : \textbf{Etienne}\\[3pt]
Titre :\\
\textbf{\MakeUppercase{Fluxional compiler : seamless shift from development productivity to performance efficiency, in the case of real-time web applications}}
\\

\noindent %
Nature : \textbf{\MakeUppercase{Doctorat}} \hfill Numéro d'ordre :  \textbf{AAAALYSEIXXXX}\\[3pt]
Ecole doctorale : \textbf{INFOMATHS}\\[3pt]
Spécialité : \textbf{Informatique}\\[3pt]
Résumé : \\
{\fontsize{11pt}{12pt}\selectfont%
La plupart des grands services web commencèrent comme de simples projets, et grossirent exponentiellement.
Dès le début, il est important de s'assurer de répondre aux besoins du marché.
Des langages comme Ruby ou Java sont devenus populaires en proposant la productivité  nécessaire pour itérer rapidement sur les retours utilisateurs.
Une application web qui répond correctement aux besoins des utilisateurs peut être adoptée de manière virale.
Cela lui demande d'être efficace pour traiter cette augmentation de trafic

Il est difficile pour une application d'être à la fois productive et efficace.
Quand l'audience devient trop importante, il deviens nécessaire de réécrire l'application vers un modèle plus efficace, tel qu'un pipeline.
Cette transformation représente un risque.
Pour éviter ce risque, cette thèse propose de maintenir conjointement les représentations productives et efficaces d'une même application.

Javascript est un langage productif avec une communauté importante.
Sa boucle évènementielle est similaire à un pipeline.
Ces deux modèles d’exécution traitent un flux de requêtes en chaînant des fonctions les unes après les autres.
Cependant, la boucle évènementielle permet une approche productive grâce à sa mémoire globale, tandis que le pipeline permet une exécution efficace du fait de sa parallélisation.

Cette thèse étudie la possibilité pour une équivalence de transformer une implémentation d'une représentation vers l'autre.
Avec cette équivalence, l'équipe de développement peut suivre les deux approches simultanément.
Elle peut itérer continuellement pour prendre en compte les avantages des deux approches.

Cette thèse présente un compilateur qui permet d'identifier un pipeline dans une application Javascript, et d'isoler chaque étape dans une fluxion.
Elle exécute une fonction pour chaque datum sur le flux.
Les fluxions sont indépendantes, et peuvent être déplacées d'une machine à l'autre pour amortir l'augmentation du trafic.
L'équipe de développement peut commencer à développer avec productivité.
Et avec la transformation, elle peut itérer pour progressivement atteindre l’efficacité.

~
}

\noindent %
Mots-clés : \textbf{Javascript, Web, Fluxions}\\[3pt]
Laboratoire (s) de recherche : \textbf{CITI}\\[3pt]
Directeur de thèse: \textbf{Stéphane \MakeUppercase{Frénot}}\\[3pt]
Président de jury : \textbf{Gaël \MakeUppercase{Thomas}} ou \textbf{Frédéric \MakeUppercase{Loulergue}~}\\[3pt]
Composition du jury :\\
\textbf{Gaël \MakeUppercase{Thomas}}\\
\textbf{Frédéric \MakeUppercase{Loulergue}}\\
\textbf{Floréal \MakeUppercase{Morandat}}\\
\textbf{Frédéric \MakeUppercase{Oblé}}\\   
\textbf{Stéphane \MakeUppercase{Frénot}}
\end{framed}

\restoregeometry
\eject